\subsection{First attempt}
The first DQN used the drop in waiting time as its reward, a 5\,s decision step and a full-lane
view. On the unseen days it reduced the mean waiting time from 485\,s (fixed-time) to 294\,s, but
the actuated controller reached 98\,s. Its reward did not see vehicles that could not enter the
road, it did not pay for a change of road, and it could leave a quiet road unserved for a long
time. These findings led to the combined reward and the strategies in Table~\ref{tab:strategies}.

\subsection{Strategy comparison}
Table~\ref{tab:ablation} and Figure~\ref{fig:ablation} compare the ten training runs with the
three baselines on the unseen days at 1.0 and 1.3 times the real demand (one route seed per day).

\begin{table}[H]
\centering\small
\begin{tabular}{lrrrrr}
\toprule
 & \multicolumn{2}{c}{Mean wait (s)} & \multicolumn{2}{c}{Throughput (veh/h)} & CO$_2$ \\
\cmidrule(lr){2-3}\cmidrule(lr){4-5}
Controller & 1.0$\times$ & 1.3$\times$ & 1.0$\times$ & 1.3$\times$ & (g/veh) \\
\midrule
\multicolumn{6}{l}{\textit{Baselines}} \\
Fixed-time & 495 & 952 & 3033 & 3459 & 481 \\
Actuated (loop sensors) & 86 & 491 & 3393 & 3968 & 300 \\
\rowcolor{bench}\textbf{Camera rule} & \textbf{66} & \textbf{453} & \textbf{3401} & \textbf{4051} & \textbf{298} \\
\midrule
\multicolumn{6}{l}{\textit{DQN training runs}} \\
R1 DQN, 5 s & 337 & 499 & 2502 & 2704 & 543 \\
R2 DQN, 2 s & 308 & 945 & 2988 & 2916 & 474 \\
\rowcolor{chosen}\textbf{R3--R5 warm start, camera view (3 seeds)} & \textbf{109} & \textbf{622} & \textbf{3371} & \textbf{3816} & \textbf{320} \\
R6 keep/switch action & 217 & 526 & 3283 & 3914 & 358 \\
R7 QR-DQN & 379 & 650 & 3104 & 3788 & 370 \\
R8 300k steps & 431 & 996 & 3097 & 3477 & 429 \\
R9 old reward & 150 & 549 & 3352 & 3889 & 358 \\
R10 full-lane view & 76 & 467 & 3408 & 4010 & 326 \\
\bottomrule
\end{tabular}
\caption{Strategy comparison on unseen days (16--18 January). Green: the camera rule, used as the
benchmark and teacher. Blue: the strategy chosen for the final test (warm start, camera view);
its row is the mean of three seeds.}
\label{tab:ablation}
\end{table}

\begin{figure}[H]
\centering
\fig{0.85\linewidth}{../results/ablation.png}
\caption{Mean waiting time of each strategy and baseline on unseen days.}
\label{fig:ablation}
\end{figure}

\paragraph{What the strategy comparison shows.}
\begin{itemize}
  \item \textbf{Warm start is the largest single gain.} Without it, DQN with the new reward (R2)
        waits 308\,s at real demand and 945\,s at 1.3$\times$ demand, close to the fixed plan.
        With the camera-rule warm start (R3--R5) the mean is 109\,s and 622\,s.
  \item \textbf{A longer view helps the learned controller.} R10 (warm start, full 236\,m view)
        waits 76\,s against 86\,s for actuated control (11\% less) and 467\,s against 491\,s at
        1.3$\times$ demand (5\% less). One camera at the stop line cannot give this view; it needs
        a second camera up the road or a higher mounting, so R10 is an upper bound and not part of
        the camera-only comparison.
  \item \textbf{The hand-written camera rule is the strongest controller in this test}
        (66\,s and 453\,s). It already beats actuated control, because it skips empty roads and
        ranks roads by PCU queue and waiting time. DQN starts from this rule but does not keep all
        of its quality; with a camera-only view the three seeds differ widely (96\,s, 87\,s and
        146\,s at real demand).
  \item \textbf{Longer training is worse, not better.} R8 (300{,}000 steps) waits 431\,s against
        109\,s for the same setup at 150{,}000 steps. The agent drifts away from the teacher's
        policy as training continues, which points to a need for early stopping or a stronger
        anchor to the teacher.
  \item \textbf{The other changes do not help in this setting}: the keep/switch action (R6),
        QR-DQN (R7) and the 5\,s step (R1). R1 has a low waiting time at 1.3$\times$ demand only
        because 37\% of its vehicles never entered the road and so have no waiting time recorded;
        its throughput is the lowest of all controllers.
  \item \textbf{Reward:} the combined reward gives a lower waiting time than the old reward at real
        demand (109\,s against 150\,s, R3--R5 against R9) but a higher one at 1.3$\times$ demand
        (622\,s against 549\,s), mostly because of one weak seed.
\end{itemize}
The full test therefore uses the camera rule and the camera-view DQN (the R3--R5 setup), which
both work with one camera at the stop line. Its result in the full test is the mean of several training seeds.


\subsection{Full test}
The full test compares the camera-view DQN (the R3--R5 setup, trained with \camseedsword{} seeds) with the
three baselines on all five scenarios, with \routeseedsword{} route seeds per day (Table~\ref{tab:summary},
Figure~\ref{fig:comparison}). Every controller in this test uses only data that one camera or one
set of loops at the junction can give.

\begin{table}[H]
\centering\scriptsize
\setlength{\tabcolsep}{3.5pt}
\begin{tabular}{llrrrrrr}
\toprule
Scenario & Controller & Wait (s) & Delay (s) & Trip (s) & Thr.\ (veh/h) & CO$_2$ (g) & Fuel (g) \\
\midrule
Unseen days & Fixed-time & 484.8 & 509.6 & 184.2 & 3024 & 477 & 151 \\
 & Actuated (loops) & 98.0 & 111.1 & 106.9 & 3368 & 298 & 94 \\
\rowcolor{bench} & Camera rule & 63.4 & 76.4 & 106.0 & 3392 & 294 & 93 \\
\rowcolor{chosen} & \textbf{DQN, camera view} & 103.7 $\pm$ 27.5 & 118.1 & 112.5 & 3361 & 310 & 98 \\
\midrule
1.3$\times$ demand & Fixed-time & 970.0 & 1001.6 & 227.6 & 3457 & 568 & 180 \\
 & Actuated (loops) & 504.4 & 522.3 & 163.2 & 3963 & 426 & 135 \\
\rowcolor{bench} & Camera rule & 507.3 & 526.5 & 175.9 & 4026 & 450 & 143 \\
\rowcolor{chosen} & \textbf{DQN, camera view} & 538.8 $\pm$ 71.1 & 559.1 & 177.7 & 3968 & 454 & 144 \\
\midrule
Unseen junction (RTO) & Fixed-time & 535.4 & 556.8 & 161.8 & 2710 & 427 & 135 \\
 & Actuated (loops) & 53.4 & 65.6 & 92.6 & 3100 & 277 & 88 \\
\rowcolor{bench} & Camera rule & 40.2 & 51.5 & 87.3 & 3101 & 262 & 83 \\
\rowcolor{chosen} & \textbf{DQN, camera view} & 48.9 $\pm$ 6.1 & 61.9 & 95.7 & 3100 & 284 & 90 \\
\midrule
Free left & Fixed-time & 503.4 & 528.3 & 181.1 & 3018 & 470 & 149 \\
 & Actuated (loops) & 98.7 & 111.7 & 106.7 & 3359 & 298 & 94 \\
\rowcolor{bench} & Camera rule & 55.8 & 68.2 & 98.8 & 3403 & 278 & 88 \\
\rowcolor{chosen} & \textbf{DQN, camera view} & 239.0 $\pm$ 36.6 & 253.5 & 120.8 & 3176 & 326 & 103 \\
\midrule
Free left, 1.3$\times$ & Fixed-time & 977.1 & 1007.9 & 217.9 & 3445 & 547 & 173 \\
 & Actuated (loops) & 512.9 & 530.5 & 157.3 & 3936 & 413 & 131 \\
\rowcolor{bench} & Camera rule & 431.8 & 450.0 & 159.2 & 4035 & 415 & 132 \\
\rowcolor{chosen} & \textbf{DQN, camera view} & 737.2 $\pm$ 36.8 & 757.1 & 181.6 & 3391 & 458 & 145 \\
\bottomrule
\end{tabular}
\caption{Full test results. Mean over 3 days $\times$ \routeseeds{} route seeds; for DQN also over \trainseeds{} training
seeds ($\pm$ is the standard deviation between training seeds). Thr.\ is throughput; CO$_2$ and
fuel are per vehicle.}
\label{tab:summary}
\end{table}

\begin{figure}[H]
\centering
\fig{\linewidth}{../results/comparison.png}
\caption{Waiting time, trip time, throughput and CO$_2$ for each controller and scenario.
Error bars show the spread between the \camseedsword{} training seeds of the DQN.}
\label{fig:comparison}
\end{figure}

\begin{figure}[H]
\centering
\fig{0.8\linewidth}{../results/learning_curve.png}
\caption{Episode return of the camera-view DQN during training. Thin lines: single training
seeds; thick line: their mean.}
\label{fig:learning}
\end{figure}

\begin{table}[H]
\centering\small
\begin{tabular}{lrrrr}
\toprule
Controller & Car & Motorbike & Bus & Truck \\
\midrule
Fixed-time & 525.2 & 485.2 & 402.2 & 526.4 \\
Actuated (loops) & 111.5 & 103.2 & 73.6 & 103.6 \\
Camera rule & 67.7 & 67.9 & 65.2 & 63.3 \\
\textbf{DQN, camera view} & 116.2 & 110.3 & 81.5 & 106.9 \\
\bottomrule
\end{tabular}
\caption{Mean waiting time (s) by vehicle class on the unseen days.}
\label{tab:class}
\end{table}

\subsection{Discussion}
\begin{itemize}
  \item \textbf{All adaptive controllers are far better than the fixed plan.} On the unseen days
        the fixed plan gives a mean wait of 485\,s; the actuated controller, the camera rule and
        the camera-view DQN give 98\,s, 63\,s and 104\,s. They also move about 11\% more vehicles
        per hour and use about 35\% less fuel per vehicle.
  \item \textbf{The camera rule is the best controller.} With only a 75\,m camera view it waits
        35\% less than the loop-sensor actuated controller on the unseen days (63\,s against
        98\,s), 25\% less on the unseen junction (40\,s against 53\,s) and 43\% less with free
        left turns. At 1.3$\times$ demand it equals the actuated controller (507\,s against 504\,s)
        and has a higher throughput. Its advantage comes from two choices: it skips roads with no
        queue, and it ranks roads by the PCU of the queue and the longest wait.
  \item \textbf{The camera-view DQN is close to the actuated controller.} On the unseen junction it
        waits 49\,s against 53\,s, and on the unseen days 104\,s against 98\,s. The result changes
        between training seeds (standard deviation 28\,s on the unseen days), so the mean of
        several seeds is reported.
  \item \textbf{The DQN does not handle a change of road layout.} With free left turns, which it
        never saw in training, its wait rises to 239\,s, while the two rules do not need training
        and stay good. A learned controller must train on the layout where it will run.
  \item \textbf{Buses and trucks are not left behind.} For every adaptive controller, buses wait
        less than cars and trucks wait about as long as cars (Table~\ref{tab:class}).
\end{itemize}

